\begin{example}[The unique map from an affine scheme]\label{ex:vi21-e01}
For $X=\Spec A$, the structure morphism is induced by the unique unital map $\mathbb Z\to A$.  On points it sends $\mathfrak p$ to $\mathfrak p\cap\mathbb Z$.
\end{example}
\begin{example}[A characteristic-zero point]\label{ex:vi21-e02}
The generic point $(0)\in\Spec\mathbb Z$ has residue field $\mathbb Q$, so its image under the structure morphism is $(0)$.
\end{example}
\begin{example}[A characteristic-$p$ point]\label{ex:vi21-e03}
The closed point $(p)\in\Spec\mathbb Z$ has residue field $\mathbb F_p$ and maps to $(p)$.
\end{example}
\begin{example}[Points of the affine line]\label{ex:vi21-e04}
A point $\mathfrak p\in\Spec k[t]$ carries the local ring $k[t]_{\mathfrak p}$ and residue field $\kappa(\mathfrak p)$.
\end{example}
\begin{example}[A rational point]\label{ex:vi21-e05}
For $a\in k$, evaluation $k[t]\to k$, $t\mapsto a$, gives a $k$-point $\Spec k\to\mathbb A^1_k$ with image $(t-a)$.
\end{example}
\begin{example}[A non-rational closed point]\label{ex:vi21-e06}
Over $k=\mathbb R$, the prime $(t^2+1)$ has residue field $\mathbb C$; it is closed but is not an $\mathbb R$-rational point.
\end{example}
\begin{example}[The generic point of an integral affine scheme]\label{ex:vi21-e07}
If $A$ is a domain, the point $(0)$ has residue field $\operatorname{Frac}(A)$ and specializes to every point of $\Spec A$.
\end{example}
\begin{example}[Specialization in $\Spec\mathbb Z$]\label{ex:vi21-e08}
The generic point $(0)$ specializes to every closed point $(p)$.
\end{example}
\begin{example}[Residue field under a localization morphism]\label{ex:vi21-e09}
For $A\to A_f$ and $\mathfrak q\in\Spec A_f$, the residue field of $\mathfrak q$ is canonically isomorphic to the residue field of its contraction in $\Spec A$.
\end{example}
\begin{example}[Residue field under a quotient morphism]\label{ex:vi21-e10}
For $A\to A/I$ and $\mathfrak q\supset I$, the corresponding point of $\Spec(A/I)$ has the same residue field as $\mathfrak q\in\Spec A$.
\end{example}
\begin{example}[A point morphism]\label{ex:vi21-e11}
For $x\in X$, the map $\Spec\kappa(x)\to X$ is canonical even when $x$ is not closed.
\end{example}
\begin{example}[Extension of residue fields]\label{ex:vi21-e12}
If $K/\kappa(x)$ is a field extension, then $\Spec K\to X$ has image $x$.
\end{example}
\begin{example}[Two field-valued points with the same image]\label{ex:vi21-e13}
Different embeddings $\kappa(x)\to K$ can define distinct morphisms $\Spec K\to X$ with the same underlying image point.
\end{example}
\begin{example}[A local map on stalks]\label{ex:vi21-e14}
If $f:X\to Y$ and $x\mapsto y$, then $f_x^\#:\OO_{Y,y}\to\OO_{X,x}$ sends nonunits to nonunits in the precise local-homomorphism sense.
\end{example}
\begin{example}[Characteristic map on an arithmetic affine line]\label{ex:vi21-e15}
A prime $\mathfrak p\subset\mathbb Z[t]$ maps to $\mathfrak p\cap\mathbb Z$, either $(0)$ or $(p)$.
\end{example}
\begin{example}[A scheme entirely in characteristic $p$]\label{ex:vi21-e16}
For an $\mathbb F_p$-scheme, every residue field has characteristic $p$, so the absolute structure morphism factors through $\Spec\mathbb F_p\to\Spec\mathbb Z$.
\end{example}
\begin{example}[One-point nonreduced scheme]\label{ex:vi21-e17}
$\Spec k[\varepsilon]/(\varepsilon^2)$ and $\Spec k$ have the same one-point topology but different local rings.  Point sets alone do not determine schemes.
\end{example}
\begin{example}[Continuous maps are not enough]\label{ex:vi21-e18}
A continuous map of underlying spaces does not determine the sheaf homomorphism, and a sheaf homomorphism without local stalk maps need not be a scheme morphism.
\end{example}
\begin{example}[Closed versus rational]\label{ex:vi21-e19}
Closedness is topological; rationality is residue-field-theoretic.  A closed point may have residue field strictly larger than the base field.
\end{example}
\begin{example}[Naturality of the absolute base]\label{ex:vi21-e20}
For every scheme morphism $f:X\to Y$, uniqueness gives $\pi_Y\circ f=\pi_X$.
\end{example}
